\section{Introduction}

% DO NOT MOVE THIS FIGURE - This is the hero image and must remain at the top of page 2
% Source: experiments/scripts/make_fig_hero_real.py (generates fig_hero_tikz.pdf)
\begin{figure*}[t!]
\centering
\includegraphics[width=\textwidth]{fig_hero_tikz.pdf}
\caption{Training leaves traces in residual networks. (a) The residual branch product $M = W_{\mathrm{out}} W_{\mathrm{in}}$ for a block with skip connection; the normalized trace concentration $s(M)$ measures identity alignment. (b) Training dynamics: score matrices $s(i,j)$ from epoch 0 to 200. ResNet (top row) develops strong identity-aligned structure; PlainNet without skip connections (bottom row) remains random. The fingerprint is learned, not an initialization artifact. (c) Lineage verification: the lineage score $\mathcal{L}$ (Eq.~\ref{eq:lineage-score}) aggregates per-block signature similarity. Related checkpoints (fine-tuned, quantized, pruned, LoRA-merged) score high; unrelated checkpoints (independent, distilled) score low, enabling clean separation.}
\label{fig:hero}
\end{figure*}

The open-weight ecosystem has transformed how neural networks are distributed and reused. When Meta released LLaMA~\citep{touvron2023llama}, the community produced instruction-tuned variants within weeks. Platforms like Hugging Face now host hundreds of thousands of model checkpoints~\citep{huggingface2024}, many of which are fine-tuned~\citep{devlin2019bert,ouyang2022rlhf}, quantized~\citep{han2016deep,dettmers2022llmint8}, pruned~\citep{han2015deep}, or merged~\citep{wortsman2022model,ilharco2023editing}. This creates a \emph{model supply chain} where downstream checkpoints inherit weights from upstream models, but the lineage is often undocumented or intentionally obscured. Provenance of AI-generated content is now treated as a first-class problem: Google's SynthID watermarks text and images at generation time~\citep{dathathri2024synthid}, yet no analogous mechanism exists for the models themselves. Model reuse is easy; reliable model provenance is not.

Existing provenance mechanisms fall short (Table~\ref{tab:comparison}). Cryptographic hashes break after any weight modification. Metadata and model cards~\citep{mitchell2019model} require honest distributors and preserved records; they can be falsified or lost. Watermarks~\citep{uchida2017embedding,adi2018turning} must be inserted before release; they cannot be applied retroactively to already-deployed checkpoints, and no surveyed scheme proved robust against fine-tuning or pruning~\citep{lukas2022sok}. Behavioral tests that compare model outputs cannot distinguish weight inheritance from imitation: two models can behave similarly without sharing weights, especially after distillation~\citep{hinton2015distilling}. Activation-based similarity measures like CKA~\citep{kornblith2019similarity} and decision-boundary fingerprints like IPGuard~\citep{cao2021ipguard} inherit this limitation. We need a passive, post-hoc signal that reads provenance from the weights themselves.

We ask: given a reference checkpoint and a suspect checkpoint with the same residual architecture (e.g., both LLaMA-2-7B derivatives), can we determine whether the suspect inherited weights from the reference, even after common transformations like fine-tuning, quantization, or pruning? This is distinct from duplicate detection (which hashes solve until the first modification), behavioral similarity (which conflates imitation with inheritance), and watermark verification (which requires forethought). The question is whether shared weight ancestry leaves detectable structure after modification.

Prior work~\citep{park2026idroppedneuralnet} observed that correctly paired projections in trained residual blocks develop trace-concentrated products and showed this structure recovers projection correspondences. We build on this observation and ask whether the checkpoint-specific component of these products persists across derived models and can support lineage verification. Our key step is removing the identity-aligned component shared across trained models and comparing centered signatures. Section~\ref{sec:exp-exists} validates the underlying phenomenon through initialization ablations, skip-free controls, and gradient-coupling interventions.

We exploit this structure for provenance verification. Correctly paired projections have high trace-concentration scores; mismatched projections do not. However, independently trained models exhibit the same generic identity-aligned structure, so trace concentration alone does not establish ancestry. We therefore center each branch product and compare checkpoint-specific signatures across models, aggregating similarities into a lineage score (Equation~\ref{eq:lineage-score}). Related checkpoints score near 1; unrelated ones score near 0.

The pairing signal generalizes: across public checkpoints in language, vision, and speech---GPT-2, BERT, LLaMA-2, Mistral, Qwen2.5, DeepSeek-R1, ViT, Whisper, ResNets---architecture-aware products recover projection correspondences with 80--100\% accuracy (Table~\ref{tab:transformer-family}). On controlled MLP, ResNet-18, and GPT-2-style language model benchmarks, checkpoints with genuine ancestry (fine-tuning, noise, pruning, quantization, LoRA merging) separate cleanly from independent runs and distilled students (AUROC = 1.000). Distillation is instructive: in our controlled GPT-2 benchmark, students achieve 79\% top-1 agreement and perplexity ratio 0.99 with their teachers yet score $\mathcal{L}{=}0.002$, since their weights carry no trace of the teacher's training (Figure~\ref{fig:hero}c). On public LLMs, LLaMA-2-7B derivatives (chat, Vicuna, CodeLlama) score related while OpenLLaMA-7B, trained from scratch, is correctly rejected.

Our contributions:
\begin{itemize}[leftmargin=*,nosep]
\item \textbf{We extend residual-product structure into a checkpoint-comparison framework.} Building on prior observations of trace concentration, we isolate the centered, checkpoint-specific component for lineage verification and validate the phenomenon through controlled interventions.
\item \textbf{We introduce a centered residual-signature test for white-box lineage verification.} Removing the identity-aligned component common across trained models, we compare checkpoint-specific geometry without training data, forward passes, or watermarks.
\item \textbf{We evaluate projection pairing and lineage verification at appropriate scopes.} Projection correspondences reach 80--100\% accuracy across nine families; controlled MLP, ResNet, and GPT-2-style language model benchmarks with known ancestry evaluate checkpoint-level lineage (AUROC$=$1.000), complemented by validation on public LLaMA-2 checkpoints.
\item \textbf{We show the signal survives substantial modification.} It persists through fine-tuning, quantization, pruning, LoRA merging, and RLHF, and tracks partial inheritance in layer grafts and model merges.
\end{itemize}
